\section{Method for Winding Distribution} \label{sec:WindingDist}
The developed approach is based on the analysis of the magnetic field within the machine, specifically examining the magnetic flux linkage, denoted as \(\gls{mag_fluxlink}\). Similar to the slot star, a phasor representation is used; however, instead of placing the phasors in the slots, they are aligned with the mechanical coil axis, as illustrated in Fig. \ref{fig:Zahnstern_Mechanisch}\footnote{The shown phasors are arranged radially, whereas they would need to be aligned axially, normal to the plane of the paper, according to physical constraints.}. Here, \(\gls{phasor_diff_angle_mech} = \frac{2\pi}{\gls{Ntooth}}\) represents the angular distance between two consecutive phasors in mechanical coordinates. The next step involves transforming the mechanical phasor star into electrical coordinates by multiplying the mechanical phasor argument \(\gls{gamma_mech}\) by the number of pole pairs \(\gls{polepair}\) to determine the electrical phasor argument \(\gls{gamma_elec}\): \(\gls{gamma_elec}_{k} = \gls{polepair} \cdot \gls{gamma_mech}_{k} \quad k = 1, \dots, \gls{Ntooth}\). Accordingly, the distance between two consecutive phasors in electrical coordinates is given by \(\gls{phasor_diff_angle_elec} = \gls{polepair} \cdot \gls{phasor_diff_angle_mech}\). In the next step, magnetic fluxes are assigned to the phasors in the electrical coordinate system. These fluxes contain information about the phases of the three-phase system as well as the polarity of the magnetic flux. The phasors are selected to form a rotating magnetic field, with each phase and polarity having an equal number of coils. Additionally, the vectorial sum of the phasors in each phase must result in the magnitude maximum :
%
\begin{equation}
    \pNorm
    {
        \phasor{\gls{mag_fluxlink}}[_{x+}^{\discretestep{k_1}}] +
        \phasor{\gls{mag_fluxlink}}[_{x+}^{\discretestep{k_2}}] -
        \phasor{\gls{mag_fluxlink}}[_{x-}^{\discretestep{k_3}}] -
        \phasor{\gls{mag_fluxlink}}[_{x-}^{\discretestep{k_4}}]
    }[2] \to \max \quad x = \gls{idx_cs_abc_a}, \gls{idx_cs_abc_b}, \gls{idx_cs_abc_c}.
    \label{equ:Bopt}
\end{equation}
%
Here, \(\phasor{\gls{mag_fluxlink}}[_{x\{+,-\}}^{\discretestep{k_i}}]\) is interpreted such that the subscript \(x\) represents the phase, the subscript \(\{+,-\}\) indicates the polarity, and determines whether addition or subtraction is used. The superscript \(\discretestep{k_i} \in [1;\gls{Ntooth}] \subseteq \mathbb{N}^{+}\) represents the number of the corresponding phasor. Therefore, it is necessary to choose the values \(k_1, \dots, k_4\) such that the magnitude condition is satisfied. In the case of the reference machine with \(\gls{Ntooth}=12\), \(\gls{polepair}=5\), and \(\gls{num_phases}=3\), this implies that the number of coils per phase and polarity is \(\frac{\gls{Ntooth}}{2\gls{num_phases}} = 2\), and the phasor selection for each phase is shifted by \(\SI{120}{\degree}\). This results in the flux distribution shown in Fig. \ref{fig:Zahnstern_Elektrisch}.
%
\begin{figure}[ht]
    \centering
    \begin{subfigure}{.49\linewidth}
        \centering
    \begin{tikzpicture}[font=\footnotesize]
    \pgfmathsetmacro{\slots}{12}
    \pgfmathsetmacro{\innerRadius}{1.3}
    \pgfmathsetmacro{\outerRadius}{2}
    \pgfmathsetmacro{\meanRadius}{0.5*(\innerRadius + \outerRadius)}
    \pgfmathsetmacro{\segmentangle}{360/\slots}
    \pgfmathsetmacro{\slotWidth}{12}
    \pgfmathsetmacro{\innerSlotwidth}{\slotWidth/\innerRadius}
    \pgfmathsetmacro{\outerSlotwidth}{\slotWidth/\outerRadius}


    \pgfmathsetmacro{\iters}{\slots-1}
    \foreach \k in {0,...,\iters}
    {

        \pgfmathsetmacro{\innerCW}{\k*\segmentangle - (\segmentangle - \innerSlotwidth)/2}
        \pgfmathsetmacro{\innerCCW}{\k*\segmentangle + (\segmentangle - \innerSlotwidth)/2}
        \pgfmathsetmacro{\outerCW}{\k*\segmentangle - (\segmentangle - \outerSlotwidth)/2}
        \pgfmathsetmacro{\outerCCW}{\k*\segmentangle + (\segmentangle - \outerSlotwidth)/2}
    
        \filldraw[fill=colIron] 
            (\innerCW:\innerRadius) arc[start angle=\innerCW, end angle=\innerCCW, radius=\innerRadius] --
            (\outerCCW:\outerRadius) arc[start angle=\outerCCW, end angle=\outerCW, radius=\outerRadius] -- 
            cycle;
        
        \pgfmathsetmacro{\number}{int(\k+1)}
        \draw[->, line width=1] (0,0) 
            -- ({\segmentangle*\k}:{\outerRadius*1.15}) 
            node[pos=1.1] {\number};
    }
    \pgfmathsetmacro{\arcpos}{0.8*\innerRadius}
    \draw[<->] (0:\arcpos) arc(0:30:\arcpos);
    \node at (15:{0.8*\arcpos}) {\scriptsize \gls{phasor_diff_angle_mech}};

    \path (0,0) circle (1.35*\outerRadius);
\end{tikzpicture}
        \caption{Mechanical flux star}
        \label{fig:Zahnstern_Mechanisch}
    \end{subfigure}
    \hfill
    \begin{subfigure}{.49\linewidth}
        \centering
    \begin{tikzpicture}[font=\footnotesize]
    \pgfmathsetmacro{\slots}{12}
    \pgfmathsetmacro{\polepairs}{5}
    \pgfmathsetmacro{\coveragefactor}{0.85}
    \pgfmathsetmacro{\outerRadius}{1.5}
    \pgfmathsetmacro{\segmentangle}{360/\slots}
    \pgfmathsetmacro{\slotangle}{\segmentangle*\coveragefactor}
    
    \def\Upos{\color{colPhaU}$\gls{mag_fluxlink}_{\gls{idx_cs_abc_a}+}$}
    \def\Uneg{\color{colPhaU}$\gls{mag_fluxlink}_{\gls{idx_cs_abc_a}-}$}
    \def\Vpos{\color{colPhaW}$\gls{mag_fluxlink}_{\gls{idx_cs_abc_c}+}$}
    \def\Vneg{\color{colPhaW}$\gls{mag_fluxlink}_{\gls{idx_cs_abc_c}-}$}
    \def\Wpos{\color{colPhaV}$\gls{mag_fluxlink}_{\gls{idx_cs_abc_b}+}$}
    \def\Wneg{\color{colPhaV}$\gls{mag_fluxlink}_{\gls{idx_cs_abc_b}-}$}

    \draw[->, line width=1] (0,0) -- ({\segmentangle*0*\polepairs}:{\outerRadius*1.1}) node[pos=1.1] {1} node[pos=1.35] {\Upos};
    \draw[->, line width=1] (0,0) -- ({\segmentangle*1*\polepairs}:{\outerRadius*1.1}) node[pos=1.1] {2} node[pos=1.35] {\Wpos};
    \draw[->, line width=1] (0,0) -- ({\segmentangle*2*\polepairs}:{\outerRadius*1.1}) node[pos=1.1] {3} node[pos=1.35] {\Wneg};
    \draw[->, line width=1] (0,0) -- ({\segmentangle*3*\polepairs}:{\outerRadius*1.1}) node[pos=1.1] {4} node[pos=1.35] {\Vneg};
    \draw[->, line width=1] (0,0) -- ({\segmentangle*4*\polepairs}:{\outerRadius*1.1}) node[pos=1.1] {5} node[pos=1.35] {\Vpos};
    \draw[->, line width=1] (0,0) -- ({\segmentangle*5*\polepairs}:{\outerRadius*1.1}) node[pos=1.1] {6} node[pos=1.35] {\Upos};
    \draw[->, line width=1] (0,0) -- ({\segmentangle*6*\polepairs}:{\outerRadius*1.1}) node[pos=1.1] {7} node[pos=1.35] {\Uneg};
    \draw[->, line width=1] (0,0) -- ({\segmentangle*7*\polepairs}:{\outerRadius*1.1}) node[pos=1.1] {8} node[pos=1.35] {\Wneg};
    \draw[->, line width=1] (0,0) -- ({\segmentangle*8*\polepairs}:{\outerRadius*1.1}) node[pos=1.1] {9} node[pos=1.35] {\Wpos};
    \draw[->, line width=1] (0,0) -- ({\segmentangle*9*\polepairs}:{\outerRadius*1.1}) node[pos=1.1] {10} node[pos=1.35] {\Vpos};
    \draw[->, line width=1] (0,0) -- ({\segmentangle*10*\polepairs}:{\outerRadius*1.1}) node[pos=1.1] {11} node[pos=1.35] {\Vneg};
    \draw[->, line width=1] (0,0) -- ({\segmentangle*11*\polepairs}:{\outerRadius*1.1}) node[pos=1.1] {12} node[pos=1.35] {\Uneg};    

    \pgfmathsetmacro{\arcpos}{0.7*\outerRadius}
    \draw[<->] (0:\arcpos) arc(0:30:\arcpos);
    \node at (15:{1.2*\arcpos}) {\scriptsize \gls{phasor_diff_angle_mech}};

    \pgfmathsetmacro{\arcpos}{0.6*\outerRadius}
    \draw[<->] (0:\arcpos) arc(0:150:\arcpos);
    \node at (75:{1.2*\arcpos}) {\scriptsize \gls{phasor_diff_angle_elec}};

    \path (0,0) circle (1.35*2); % Anpassen damit mechanische Mitte = ELektrischer Mitte 
\end{tikzpicture}
        \caption{Electrical flux star}
        \label{fig:Zahnstern_Elektrisch}
    \end{subfigure}
    \caption{Representation of the phasor star}
    \label{fig:Zahnstern}
\end{figure}
%
The result in Fig. \ref{fig:Zahnstern_Elektrisch} can be interpreted such that the winding, depending on the polarity, is either wound in a left-handed or right-handed direction. If one arbitrarily defines that a positive polarity generates magnetic flux upwards and a negative polarity generates magnetic flux downwards, the winding can be characterized using the ampère's right-hand grip rule. This method results in the winding distribution as shown in Fig. \ref{fig:Wicklungsverteilung}.
%
\begin{figure}[ht]
    \centering
    \begin{tikzpicture}[
    mag/.style={draw=black},
    Nmag/.style={top color=colMagnetNorth, bottom color=colMagnetSouth, mag},
    Smag/.style={top color=colMagnetSouth, bottom color=colMagnetNorth, mag},
    flux/.style={-{Computer Modern Rightarrow[]}, line width = 2.5},
    winding/.style={postaction={decorate}, decoration={markings, mark=between positions 0.175 and 0.825 step 0.325 with {\node {#1};}}},
    wpos/.style={winding={\footnotesize$\otimes$}},
    wneg/.style={winding={\footnotesize$\odot$}},
]
    \pgfmathsetmacro{\factor}{3.25e-2}
    \pgfmathsetmacro{\N}{12}
    \pgfmathsetmacro{\p}{5}
    \pgfmathsetmacro{\poles}{2*\p}
    \pgfmathsetmacro{\ri}{60.2*\factor}
    \pgfmathsetmacro{\ro}{98.8*\factor}
    \pgfmathsetmacro{\rm}{0.5*(\ri+\ro)}
    \pgfmathsetmacro{\mcf}{0.84}
    \pgfmathsetmacro{\hpm}{8.2*\factor}
    \pgfmathsetmacro{\hyoke}{7.7*\factor}
    \pgfmathsetmacro{\g}{2*\factor}

    \pgfmathsetmacro{\hshoe}{5*\factor}
    \pgfmathsetmacro{\hcore}{25*\factor}
    \pgfmathsetmacro{\wcoil}{8*\factor}
    \pgfmathsetmacro{\wslot}{2*\factor}
    \pgfmathsetmacro{\wsec}{2*pi/\N*\rm }
    \pgfmathsetmacro{\wcore}{2*pi/\N*\rm - (2*\wcoil + \wslot)}
    \pgfmathsetmacro{\wshoe}{2*pi/\N*\rm - \wslot}

    \pgfmathsetmacro{\magSecWidth}{pi/\p*\rm}
    \pgfmathsetmacro{\wmag}{\magSecWidth*\mcf}
    \pgfmathsetmacro{\wmagair}{\magSecWidth*(1-\mcf)}

    \pgfmathsetmacro{\axlen}{\hcore/2 + \hshoe + \g + \hpm + \hyoke}

    \pgfmathsetmacro{\curveangle}{12.5}

    \pgfmathsetmacro{\fluxlen}{\axlen}

    \newcommand{\phaseArray}{{1},{2},{-2},{-3},{3},{1},{-1},{-2},{2},{3},{-3},{-1}}
    % Poles
    \foreach \x in {1,...,\poles}
    {
        \pgfmathsetmacro{\xleft}{-\wmag/2 + \magSecWidth*(\x-5.5)}
        \ifthenelse{\isodd{\x}}
        {
            \draw[Nmag] (\xleft, {(\hcore/2 + \hshoe + \g + \hpm)}) rectangle ++(\wmag, {-\hpm});
            \draw[Nmag] (\xleft, {-(\hcore/2 + \hshoe + \g + \hpm)}) rectangle ++(\wmag, {+\hpm});
        }
        {
            \draw[Smag] (\xleft, {(\hcore/2 + \hshoe + \g + \hpm)}) rectangle ++(\wmag, {-\hpm});
            \draw[Smag] (\xleft, {-(\hcore/2 + \hshoe + \g + \hpm)}) rectangle ++(\wmag, {+\hpm});
        }
    }
     
    % Rotoryoke
    \draw[fill=colIron] (-6*\wsec, {-(\hcore/2 + \hshoe + \g + \hpm)})
    rectangle ++(12*\wsec, -\hyoke)
    ;
    \draw[fill=colIron] (-6*\wsec, {(\hcore/2 + \hshoe + \g + \hpm)})
    rectangle ++(12*\wsec, \hyoke)
    ;

    % Statorteeth
    \foreach \pha [count=\x] in \phaseArray
    {
        \draw[fill=colIron] ({-\wcore/2 + \wsec*(\x-6.5)}, -\hcore/2)
            -- ++(0, \hcore) coordinate(P_NE_L)
            -- ++(-\wcoil, 0)
            -- ++(0, \hshoe)
            -- ++(\wshoe, 0)
            -- ++(0, -\hshoe) coordinate(P_NE_R)
            -- ++(-\wcoil, 0)
            -- ++(0, -\hcore) coordinate(P_SW_R)
            -- ++(\wcoil, 0)
            -- ++(0, -\hshoe)
            -- ++(-\wshoe, 0)
            -- ++(0, \hshoe) coordinate(P_SW_L)
            -- cycle
            ;

        \coordinate(center) at ($(P_SW_L)!0.5!(P_NE_R)$);
        \coordinate(lc) at ($(P_NE_L)!0.6!(P_SW_L)$);
        \coordinate(rc) at ($(P_NE_R)!0.4!(P_SW_R)$);
            
        \ifnum \pha = -3
            \path[wpos, colPhaW] (lc|-P_NE_L) -- (lc|-P_SW_L);
            \path[wneg, colPhaW] (rc|-P_NE_R) -- (rc|-P_SW_R);
            \draw[flux, colPhaW] ($(0,\fluxlen/2) + (center)$) -- ++(0,-\fluxlen);
        \else\ifnum \pha = -2
            \path[wpos, colPhaV] (lc|-P_NE_L) -- (lc|-P_SW_L);
            \path[wneg, colPhaV] (rc|-P_NE_R) -- (rc|-P_SW_R);
            \draw[flux, colPhaV] ($(0,\fluxlen/2) + (center)$) -- ++(0,-\fluxlen);
        \else\ifnum \pha = -1
            \path[wpos, colPhaU] (lc|-P_NE_L) -- (lc|-P_SW_L);
            \path[wneg, colPhaU] (rc|-P_NE_R) -- (rc|-P_SW_R);
            \draw[flux, colPhaU] ($(0,\fluxlen/2) + (center)$) -- ++(0,-\fluxlen);
        \else\ifnum \pha = 1
            \path[wneg, colPhaU] (lc|-P_NE_L) -- (lc|-P_SW_L);
            \path[wpos, colPhaU] (rc|-P_NE_R) -- (rc|-P_SW_R);
            \draw[flux, colPhaU] ($(0,-\fluxlen/2) + (center)$) -- ++(0,\fluxlen);
        \else\ifnum \pha = 2
            \path[wneg, colPhaV] (lc|-P_NE_L) -- (lc|-P_SW_L);
            \path[wpos, colPhaV] (rc|-P_NE_R) -- (rc|-P_SW_R);
            \draw[flux, colPhaV] ($(0,-\fluxlen/2) + (center)$) -- ++(0,\fluxlen);
        \else\ifnum \pha = 3
            \path[wneg, colPhaW] (lc|-P_NE_L) -- (lc|-P_SW_L);
            \path[wpos, colPhaW] (rc|-P_NE_R) -- (rc|-P_SW_R);
            \draw[flux, colPhaW] ($(0,-\fluxlen/2) + (center)$) -- ++(0,\fluxlen);
        \fi\fi\fi\fi\fi\fi

        \node[fill=white, inner sep=2pt] at (center) {\x};
    }
\end{tikzpicture}


    \caption{Winding distribution}
    \label{fig:Wicklungsverteilung}
\end{figure}
%
The developed method is limited to the application of concentrated windings. However, for this use case, it provides a straightforward method for determining the winding distribution. An additional advantage is that the \(\gls{idx_cs_ab_a}\)-axis of the stator-fixed coordinate system can be directly determined by summing the phasors of the first phase. This is of paramount importance for the analysis of the electric machine.
